\documentclass[11pt,openany]{report}
\usepackage[letterpaper,margin=0.82in,headheight=14pt]{geometry}
\usepackage[T1]{fontenc}
\usepackage[utf8]{inputenc}
\usepackage{lmodern}
\usepackage{microtype}
\usepackage{amsmath,amssymb,amsthm,mathtools}
\usepackage{booktabs,longtable,tabularx,array,multirow}
\usepackage{graphicx,float,placeins}
\usepackage[dvipsnames]{xcolor}
\usepackage{enumitem}
\usepackage{fancyhdr}
\usepackage{listings}
\usepackage{url}
\usepackage[hidelinks]{hyperref}
\usepackage[nameinlink,noabbrev]{cleveref}
\usepackage{needspace}
\allowdisplaybreaks

\hypersetup{
  pdftitle={Projector-Generated Special Functions IV},
  pdfauthor={Komorebi P.},
  pdfsubject={Squarefree depth-time kernels, Gamma dephasing, Mobius resolvents, and signed Pade packet models}
}

\definecolor{Ink}{HTML}{17202A}
\definecolor{Blue}{HTML}{2E86AB}
\definecolor{Teal}{HTML}{2A9D8F}
\definecolor{Orange}{HTML}{F4A261}
\definecolor{Magenta}{HTML}{A23E8C}
\definecolor{Red}{HTML}{E76F51}
\definecolor{Paper}{HTML}{FAFAF7}
\definecolor{Muted}{HTML}{5D6D7E}
\definecolor{CodeBG}{HTML}{F3F5F5}

\pagecolor{Paper}
\color{Ink}
\graphicspath{{figures/}}
\setlength{\parindent}{0pt}
\setlength{\parskip}{0.52em}
\raggedbottom

\pagestyle{fancy}
\fancyhf{}
\fancyhead[L]{\small Projector-Generated Special Functions IV}
\fancyhead[R]{\small Komorebi P.}
\fancyfoot[C]{\thepage}
\renewcommand{\headrulewidth}{0.4pt}

\newtheorem{theorem}{Theorem}[chapter]
\newtheorem{proposition}[theorem]{Proposition}
\newtheorem{lemma}[theorem]{Lemma}
\newtheorem{corollary}[theorem]{Corollary}
\newtheorem{conjecture}[theorem]{Conjecture}
\theoremstyle{definition}
\newtheorem{definition}[theorem]{Definition}
\newtheorem{example}[theorem]{Worked example}
\newtheorem{algorithm}[theorem]{Algorithm}
\theoremstyle{remark}
\newtheorem{remark}[theorem]{Remark}
\newtheorem{boundary}[theorem]{Boundary statement}

\newcommand{\C}{\mathbb C}
\newcommand{\R}{\mathbb R}
\newcommand{\N}{\mathbb N}
\newcommand{\e}{\mathrm e}
\newcommand{\ii}{\mathrm i}
\newcommand{\Kcal}{\mathcal K}
\newcommand{\Ccal}{\mathcal C}
\newcommand{\Mcal}{\mathcal M}
\newcommand{\Dcal}{\mathcal D}
\newcommand{\norm}[1]{\left\lVert#1\right\rVert}
\newcommand{\abs}[1]{\left\lvert#1\right\rvert}
\newcommand{\dd}{\,\mathrm d}
\newcommand{\diag}{\operatorname{diag}}
\newcommand{\sech}{\operatorname{sech}}
\newcommand{\Expect}{\mathbb E}
\newcommand{\Prob}{\mathbb P}
\newcommand{\ket}[1]{\left|#1\right\rangle}
\newcommand{\bra}[1]{\left\langle#1\right|}

\newcolumntype{L}[1]{>{\raggedright\arraybackslash}p{#1}}

\lstdefinestyle{pythonstyle}{
  language=Python,
  basicstyle=\ttfamily\scriptsize,
  backgroundcolor=\color{CodeBG},
  frame=single,
  rulecolor=\color{Muted!45},
  breaklines=true,
  breakatwhitespace=false,
  showstringspaces=false,
  numbers=left,
  numberstyle=\tiny\color{Muted},
  numbersep=7pt,
  tabsize=4,
  columns=fullflexible,
  keepspaces=true
}
\lstset{style=pythonstyle}

\newcommand{\FigureWide}[3]{%
\Needspace{0.53\textheight}
\begin{figure}[H]
\centering
\includegraphics[width=0.92\textwidth,height=0.44\textheight,keepaspectratio]{#1.pdf}
\caption{#2 #3}
\label{fig:#1}
\end{figure}
}

\newcommand{\InlinePart}[1]{%
\FloatBarrier
\Needspace{10\baselineskip}
\refstepcounter{part}%
\addcontentsline{toc}{part}{\protect\numberline{\thepart}#1}%
\vspace{1.4em}
{\color{Blue}\Large\bfseries Part \thepart\par}
\vspace{0.2em}
{\color{Blue}\LARGE\bfseries #1\par}
\vspace{0.9em}
}

\newcommand{\FlowChapter}[1]{%
\FloatBarrier
\Needspace{14\baselineskip}
\refstepcounter{chapter}%
\setcounter{section}{0}%
\setcounter{subsection}{0}%
\addcontentsline{toc}{chapter}{\protect\numberline{\thechapter}#1}%
\markboth{#1}{#1}%
\vspace{1.3em}
{\color{Muted}\large\bfseries Chapter \thechapter\par}
\vspace{0.15em}
{\LARGE\bfseries #1\par}
\vspace{0.9em}
}

\title{\textbf{Projector-Generated Special Functions IV}\\[0.25em]
\Large Squarefree Depth--Time Kernels, Universal Gamma Dephasing, M\"obius Resolvents, and Signed Pad\'e Packet Models}
\author{Komorebi P.}
\date{}

\begin{document}
\hypersetup{pageanchor=false}
\pagenumbering{roman}
\maketitle

\begin{abstract}
This paper develops the first joint scale-time special-function system generated by the projector-resolved calculus. The signed seed is the M\"obius transfer
\[
C_{a,\beta,\chi}(k,t)
=
\sum_{n\ge1}
\frac{\mu(n)\chi(n)}{n^{a+\ii t}}
\left(1-\frac1{n^\beta}\right)^k.
\]
Its long-time cross-correlation removes the arithmetic signs on the diagonal and produces a positive depth-time kernel. The kernel is proved to be positive definite and to admit an exact finite-difference coordinate system built from squarefree Dirichlet series. Its large-depth behavior is universal. After one carrier phase is removed, the equal-depth correlation converges to
\[
\frac{\Gamma(\alpha+\ii\tau/\beta)}{\Gamma(\alpha)},
\qquad
\alpha=\frac{2a-1}{\beta}.
\]
For the principal parameters $a=\beta=2$, the squared dephasing envelope has the elementary closed form
\[
\frac{1+\tau^2}{\cosh(\pi\tau/2)}.
\]

The paper also proves a universal overlap law between different packet depths. If $k/\ell$ tends to a positive ratio, then the normalized overlap tends to
\[
\sech^\alpha\!\left(\frac12\log\frac{k}{\ell}\right).
\]
The tilted logarithmic frequencies have explicit limiting cumulants governed by the polygamma functions. Cross-kernels between distinct Dirichlet-character sectors decay at the unconditional rate $m^{-1/(2\beta)}$ after normalization.

The signed side is encoded by a M\"obius resolvent. It satisfies an exact parameter-shift equation and a Borel--Riesz bridge. Whenever a finite Hankel determinant is nonzero, a unique Pad\'e packet model preserves the first $2m$ signed moments after $m$ packets. Unlike the positive Jacobi models of the preceding paper, signed packet nodes need not remain in the physical interval. This failure is exhibited explicitly.

The positive kernel is the long-time autocorrelation of the signed process. The Riemann hypothesis remains equivalent to the statement that the distinguished recurrence time $t=0$ is not polynomially more coherent than the average packet power. This paper proves the complete unconditional depth-time and finite-reduction framework. It does not prove the remaining M\"obius coherence bound.

All figures were generated by the accompanying Python program. The release contains 28 vector figures, matching PNG files, exact tables, 18 built-in checks, 18 independent tests, and the complete Python source.
\end{abstract}

\chapter*{Status and claim boundary}
\addcontentsline{toc}{chapter}{Status and claim boundary}
The paper proves the depth-time Gram representation, the exact time-average identity, the squarefree $L$-ratio coordinate formula, the uniform Gamma dephasing limit on compact time intervals, the principal closed-form envelope, the log-frequency cumulant limit, the cross-depth overlap law, and the character cross-kernel decay theorem. It also proves the analytic signed-resolvent representation, its parameter-shift identity, the Borel--Riesz bridge, and the finite Pad\'e packet theorem under a nonvanishing Hankel determinant.

The positive kernel theory is unconditional. The finite signed reduction is conditional only on the explicitly stated nonbreakdown determinant. None of these results proves RH. The remaining assertion is the subpolynomial growth of the principal time-zero M\"obius coherence. Historical priority and full references are reserved for the consolidated end-of-series bibliography.

The logarithmic quantum interpretation is compatible with the external 2026 construction in which the number basis carries the Hamiltonian $H\ket n=\log n\ket n$ and the phase $n^{-\ii t}$. The present work uses that coordinate but studies an autocorrelation and scale filter that are distinct from the direct zeta observable.

\tableofcontents
\listoffigures
\clearpage
\pagenumbering{arabic}
\hypersetup{pageanchor=true}

\InlinePart{The signed transfer and its positive Gram lift}

\FlowChapter{One signed seed and two exact evolutions}

Fix
\[
a>1,
\qquad
\beta>0,
\]
and let $\chi$ be a Dirichlet character. Set
\[
y_n=1-\frac1{n^\beta}.
\]
The signed transfer is
\[
\boxed{
C_{a,\beta,\chi}(k,t)
=
\sum_{n\ge1}
\mu(n)\chi(n)n^{-a-\ii t}y_n^k.
}
\]
Absolute convergence follows from $a>1$. The variable $k$ is a discrete Euclidean depth. The variable $t$ is logarithmic time.

The shift in $k$ gives the exact recursion
\[
\boxed{
C_{a,\beta,\chi}(k+1,t)
=
C_{a,\beta,\chi}(k,t)
-
C_{a+\beta,\beta,\chi}(k,t).
}
\]
Indeed,
\[
y_n^{k+1}
=
y_n^k-n^{-\beta}y_n^k.
\]

Expanding the depth filter produces the finite-difference form
\[
\boxed{
C_{a,\beta,\chi}(k,t)
=
\sum_{j=0}^{k}
(-1)^j\binom{k}{j}
\frac1{L(a+\beta j+\ii t,\chi)}.
}
\]
The M\"obius packet is therefore a weighted finite-difference orbit of reciprocal $L$-values.

\FigureWide{01_theory_pipeline}{The mathematical pipeline.}{The signed transfer is lifted by time averaging into a positive kernel. The positive kernel has a universal Gamma limit. The signed resolvent and its finite Pad\'e reductions preserve the coherent information that the positive average removes.}

\begin{example}[Principal Bernoulli--zeta orbit]
Take $a=\beta=2$ and the principal character. Then
\[
C_{2,2}(k,0)
=
\sum_{j=0}^{k}
(-1)^j\binom{k}{j}\frac1{\zeta(2j+2)}.
\]
The first values are
\[
\begin{aligned}
C_{2,2}(0,0)&=0.607927101854\ldots,\\
C_{2,2}(1,0)&=-0.316011301068\ldots,\\
C_{2,2}(2,0)&=-0.256997111725\ldots.
\end{aligned}
\]
These are the B\'aez-Duarte coefficients in the normalization used throughout the series.
\end{example}

\FlowChapter{The depth--time kernel}

For two characters $\chi$ and $\psi$, define
\[
\boxed{
\Kcal_{\chi,\psi}(k,\ell;\tau)
=
\sum_{n\ge1}
\mu(n)^2\chi(n)\overline{\psi(n)}
 n^{-2a-\ii\tau}
y_n^{k+\ell}.
}
\]
When $\chi=\psi$, the weights are nonnegative.

\begin{theorem}[Exact time-average identity]
For $a>1$,
\[
\boxed{
\lim_{T\to\infty}
\frac1{2T}
\int_{-T}^{T}
C_{a,\beta,\chi}(k,t+\tau)
\overline{C_{a,\beta,\psi}(\ell,t)}\dd t
=
\Kcal_{\chi,\psi}(k,\ell;\tau).
}
\]
\end{theorem}

\begin{proof}
Expand the product into the absolutely convergent double series. The time average of the factor
\[
\e^{-\ii t(\log n-\log m)}
\]
is one when $n=m$ and tends to zero when $n\ne m$. Only the diagonal terms remain.
\end{proof}

\begin{theorem}[Positive-definite packet kernel]
For one fixed character $\chi$, the function
\[
((k,t),(\ell,s))
\longmapsto
\Kcal_{\chi,\chi}(k,\ell;t-s)
\]
is positive definite.
\end{theorem}

\begin{proof}
Define the feature vector
\[
v_{k,t}(n)
=
\frac{\mu(n)\abs{\chi(n)}}{n^a}
y_n^k n^{-\ii t}.
\]
Then
\[
\Kcal_{\chi,\chi}(k,\ell;t-s)
=
\sum_n
v_{k,t}(n)\overline{v_{\ell,s}(n)}.
\]
Every finite Gram sum is therefore nonnegative.
\end{proof}

\FigureWide{02_feature_map}{The exact feature map.}{Depth controls the active scale. Logarithmic time rotates each number-basis packet at frequency $\log n$. Their Gram contraction is the depth-time kernel.}

\FigureWide{03_kernel_gram}{A finite depth-time Gram matrix.}{The matrix is Hermitian positive semidefinite. The plotted eigenvalues provide a numerical certificate for the displayed packet set.}

\FlowChapter{Squarefree Dirichlet coordinates}

Let
\[
\xi=\chi\overline\psi
\]
and define
\[
D_\xi(s)
=
\sum_{n\ge1}
\frac{\mu(n)^2\xi(n)}{n^s}.
\]
For $\Re s>1$,
\[
\boxed{
D_\xi(s)
=
\prod_p\left(1+\frac{\xi(p)}{p^s}\right)
=
\frac{L(s,\xi)}{L(2s,\xi^2)}.
}
\]

\begin{theorem}[Finite-difference coordinate formula]
Let $m=k+\ell$. Then
\[
\boxed{
\Kcal_{\chi,\psi}(k,\ell;\tau)
=
\sum_{j=0}^{m}
(-1)^j\binom{m}{j}
D_\xi(2a+\beta j+\ii\tau).
}
\]
\end{theorem}

\begin{proof}
Expand
\[
y_n^m
=
\left(1-n^{-\beta}\right)^m
=
\sum_{j=0}^{m}
(-1)^j\binom{m}{j}n^{-\beta j}
\]
and interchange the finite sum with the absolutely convergent Dirichlet series.
\end{proof}

\FigureWide{04_l_ratio_coordinates}{Two exact coordinates for one kernel.}{The direct squarefree sum and the finite difference of zeta ratios agree to floating-point roundoff in the worked calculation.}

\FigureWide{05_time_average}{Finite-time extraction of the diagonal kernel.}{Increasing the observation window suppresses off-diagonal logarithmic frequencies and approaches the exact positive packet correlation.}

\InlinePart{Universal depth-time asymptotics}

\FlowChapter{A complex Beta asymptotic}

Assume $\chi=\psi$ has conductor $q$. Then
\[
\abs{\chi(n)}^2=\mathbf 1_{(n,q)=1}.
\]
Let
\[
A_q(X)
=
\sum_{\substack{n\le X\\(n,q)=1}}
\mu(n)^2.
\]
The squarefree counting law gives
\[
A_q(X)
=
\delta_qX+O_q(X^{1/2}),
\]
where
\[
\delta_q
=
\frac1{\zeta(2)}
\prod_{p\mid q}\frac{p}{p+1}.
\]

Set
\[
\alpha=\frac{2a-1}{\beta}
\]
and
\[
\lambda(\tau)
=
\alpha+\frac{\ii\tau}{\beta}.
\]

\begin{theorem}[Uniform complex packet asymptotic]
For every fixed $T>0$, uniformly for $\abs\tau\le T$,
\[
\boxed{
\Kcal_{\chi,\chi}(k,\ell;\tau)
=
\frac{\delta_q}{\beta}
B\!\left(\lambda(\tau),k+\ell+1\right)
+O_{a,\beta,q,T}
\left(
(k+\ell)^{-(2a-1/2)/\beta}
\right).
}
\]
Consequently,
\[
\boxed{
\Kcal_{\chi,\chi}(k,\ell;\tau)
\sim
\frac{\delta_q}{\beta}
\Gamma\!\left(\alpha+\frac{\ii\tau}{\beta}\right)
(k+\ell)^{-\alpha-\ii\tau/\beta}.
}
\]
\end{theorem}

\begin{proof}
Apply partial summation to the function
\[
f_m(x)
=
x^{-2a-\ii\tau}
\left(1-x^{-\beta}\right)^m,
\qquad
m=k+\ell.
\]
The main term is
\[
\delta_q\int_1^\infty f_m(x)\dd x.
\]
With $u=x^{-\beta}$, this becomes
\[
\frac{\delta_q}{\beta}
\int_0^1
u^{\lambda(\tau)-1}(1-u)^m\dd u
=
\frac{\delta_q}{\beta}B(\lambda(\tau),m+1).
\]
The square-root error in $A_q(X)$, together with one derivative of $f_m$, contributes
\[
O(m^{-(2a-1/2)/\beta})
\]
uniformly on compact $\tau$ intervals. The Gamma asymptotic follows from the standard Beta quotient.
\end{proof}

\FlowChapter{The Gamma dephasing function}

Define
\[
\boxed{
\Dcal_{\alpha,\beta}(\tau)
=
\frac{\Gamma(\alpha+\ii\tau/\beta)}{\Gamma(\alpha)}.
}
\]
This is the characteristic function of $\beta^{-1}\log U$ when
\[
U\sim\operatorname{Gamma}(\alpha,1).
\]
It is therefore positive definite. Since the Gamma function has no zeros,
\[
\boxed{
\Dcal_{\alpha,\beta}(\tau)\ne0
\quad
\text{for every real }\tau.
}
\]

\begin{corollary}[Equal-depth dephasing]
Let
\[
A_k(\tau)
=
\frac{\Kcal_{\chi,\chi}(k,k;\tau)}
{\Kcal_{\chi,\chi}(k,k;0)}.
\]
Then
\[
\boxed{
(2k)^{\ii\tau/\beta}A_k(\tau)
\longrightarrow
\Dcal_{\alpha,\beta}(\tau)
}
\]
uniformly on compact real $\tau$ intervals.
\end{corollary}

\FigureWide{06_gamma_trajectory}{The limiting dephasing trajectory.}{The Gamma ratio traces a zero-free curve in the complex plane. The marked point is the normalized value at zero lag.}

\FigureWide{07_gamma_convergence}{Convergence to the Gamma ratio.}{The dephased discrete correlations approach the same real and imaginary limit as the packet depth grows.}

\begin{example}[Closed principal envelope]
For
\[
a=\beta=2,
\qquad
\alpha=\frac32,
\]
we have
\[
\Dcal_{3/2,2}(\tau)
=
\frac{\Gamma(3/2+\ii\tau/2)}{\Gamma(3/2)}.
\]
Using
\[
\abs{\Gamma(1/2+\ii y)}^2
=
\frac{\pi}{\cosh(\pi y)},
\]
one obtains
\[
\boxed{
\abs{\Dcal_{3/2,2}(\tau)}^2
=
\frac{1+\tau^2}{\cosh(\pi\tau/2)}.
}
\]
\end{example}

\FigureWide{08_principal_modulus}{The principal closed-form dephasing envelope.}{The Gamma-ratio modulus and the elementary hyperbolic-cosine expression coincide throughout the displayed range.}

The phase derivative also has a closed law:
\[
\boxed{
\frac{\dd}{\dd\tau}
\arg\Dcal_{\alpha,\beta}(\tau)
=
\frac1\beta
\Re\psi\!\left(\alpha+\frac{\ii\tau}{\beta}\right).
}
\]

\FigureWide{09_phase_digamma}{The phase and its derivative.}{The digamma function supplies the exact group-delay law of the limiting packet correlation.}

\begin{table}[H]
\centering
\caption{Selected values of the limiting Gamma dephasing function for $a=\beta=2$.}
\begin{tabular}{rrrr}
\toprule $\tau$ & Real part & Imaginary part & Squared modulus \\ \midrule
-4 & 0.18721515 & -0.16865147 & 0.06349283 \\
-2 & 0.64917367 & -0.09953509 & 0.43133369 \\
-1 & 0.89225332 & -0.03094590 & 0.79707363 \\
0 & 1.00000000 & 0.00000000 & 1.00000000 \\
1 & 0.89225332 & 0.03094590 & 0.79707363 \\
2 & 0.64917367 & 0.09953509 & 0.43133369 \\
4 & 0.18721515 & 0.16865147 & 0.06349283 \\
\bottomrule
\end{tabular}
\end{table}


\FlowChapter{Log-frequency cumulants}

Normalize the depth-$m$ positive weights by
\[
\Prob_m(n)
=
\frac{
\mu(n)^2\abs{\chi(n)}^2
n^{-2a}y_n^m
}{
\Kcal_{\chi,\chi}(m,0;0)
}.
\]
Define the centered logarithmic frequency
\[
Z_m
=
\log n-\frac1\beta\log m.
\]
The preceding characteristic-function limit gives
\[
Z_m
\Longrightarrow
-\frac1\beta\log U,
\qquad
U\sim\operatorname{Gamma}(\alpha,1).
\]

\begin{corollary}[Polygamma cumulants]
For every fixed order for which moment convergence is taken,
\[
\boxed{
\kappa_r(Z_m)
\longrightarrow
\frac{(-1)^r}{\beta^r}
\psi_{r-1}(\alpha).
}
\]
In particular,
\[
\boxed{
\Expect Z_m
\longrightarrow
-\frac{\psi(\alpha)}{\beta}
}
\]
and
\[
\boxed{
\operatorname{Var}(Z_m)
\longrightarrow
\frac{\psi_1(\alpha)}{\beta^2}.
}
\]
\end{corollary}

\FigureWide{10_log_frequency_profile}{The centered logarithmic spectrum.}{The finite squarefree packet histogram approaches the log-Gamma density predicted by the dephasing function.}

\FigureWide{11_log_cumulants}{Convergence of the first two cumulants.}{The mean and variance settle toward explicit digamma and trigamma constants.}

\FlowChapter{The cross-depth overlap law}

At zero time lag, define
\[
\mathcal O(k,\ell)
=
\frac{
\Kcal_{\chi,\chi}(k,\ell;0)
}{
\sqrt{
\Kcal_{\chi,\chi}(k,k;0)
\Kcal_{\chi,\chi}(\ell,\ell;0)
}
}.
\]

\begin{theorem}[Universal scale overlap]
If $k,\ell\to\infty$ and $k/\ell\to r\in(0,\infty)$, then
\[
\boxed{
\mathcal O(k,\ell)
\longrightarrow
\left(
\frac{2\sqrt r}{1+r}
\right)^\alpha
=
\sech^\alpha\!\left(\frac12\log r\right).
}
\]
\end{theorem}

\begin{proof}
The zero-lag kernel is regularly varying:
\[
\Kcal(k,\ell;0)
\sim
A(k+\ell)^{-\alpha}.
\]
The two diagonal norms satisfy
\[
\Kcal(k,k;0)
\sim
A(2k)^{-\alpha}
\]
and the analogous formula for $\ell$. Taking the ratio gives the result.
\end{proof}

\begin{example}[A four-to-one depth ratio]
For $a=\beta=2$, we have $\alpha=3/2$. When $k/\ell=4$,
\[
\mathcal O(k,\ell)
\longrightarrow
\left(\frac45\right)^{3/2}
=
0.7155417528\ldots.
\]
The finite calculation at depths $k=40000$ and $\ell=10000$ gives a value within two percent of this limit.
\end{example}

\FigureWide{12_depth_overlap_limit}{The exact limiting overlap profile.}{Packets at comparable depths retain substantial overlap. Packets separated multiplicatively in depth become asymptotically orthogonal.}

\FigureWide{13_depth_overlap_convergence}{Finite-depth convergence.}{Several absolute scales approach the same ratio-dependent law.}

\FigureWide{14_depth_overlap_heatmap}{The geometry on logarithmic depth.}{The overlap depends asymptotically only on the difference of log depths, which produces a translation-invariant ridge in scale coordinates.}

\InlinePart{Character packets and arithmetic decoherence}

\FlowChapter{Cross-character depth-time kernels}

For two characters modulo the same conductor, the cross-kernel has the exact time-average interpretation already proved. If
\[
\xi=\chi\overline\psi
\]
is nonprincipal, its squarefree partial sums satisfy
\[
\sum_{n\le X}\mu(n)^2\xi(n)
=O_q(X^{1/2}).
\]
This follows from
\[
\mu(n)^2
=
\sum_{d^2\mid n}\mu(d)
\]
and the bounded complete-period sums of a nonprincipal character.

\begin{theorem}[Character depth-time decoherence]
If $\chi\ne\psi$, then uniformly for $\abs\tau\le T$,
\[
\boxed{
\Kcal_{\chi,\psi}(k,\ell;\tau)
=
O_{a,\beta,q,T}
\left(
(k+\ell)^{-(2a-1/2)/\beta}
\right).
}
\]
After diagonal normalization,
\[
\boxed{
\frac{
\Kcal_{\chi,\psi}(k,\ell;\tau)
}{
\sqrt{
\Kcal_{\chi,\chi}(k,k;0)
\Kcal_{\psi,\psi}(\ell,\ell;0)
}
}
=
O\left((k+\ell)^{-1/(2\beta)}\right)
}
\]
whenever the two depths remain comparable.
\end{theorem}

\begin{proof}
Use partial summation with the nonprincipal square-root bound instead of the principal density term. The same scaled derivative estimate used in the complex Beta theorem yields the stated power. Dividing by the two diagonal norms removes $(k+\ell)^{-\alpha}$ and leaves the factor $(k+\ell)^{-1/(2\beta)}$.
\end{proof}

For $a=\beta=2$, the unconditional normalized decay rate is
\[
(k+\ell)^{-1/4}.
\]

\FigureWide{15_character_gram}{Modulo-five character Gram matrices.}{The diagonal remains normalized to one while the off-diagonal sectors weaken with increasing depth.}

\FigureWide{16_character_decoherence}{The largest off-diagonal character correlation.}{The finite values decrease in the direction predicted by the unconditional quarter-power law.}

\InlinePart{The signed M\"obius function family}

\FlowChapter{The M\"obius resolvent}

Define the signed resolvent
\[
\boxed{
\Ccal_{a,\beta,\chi}(z,t)
=
\sum_{k=0}^{\infty}
C_{a,\beta,\chi}(k,t)z^k.
}
\]
For $\abs z<1$, geometric summation gives
\[
\boxed{
\Ccal_{a,\beta,\chi}(z,t)
=
\sum_{n\ge1}
\frac{
\mu(n)\chi(n)n^{-a-\ii t}
}{
1-z(1-n^{-\beta})
}.
}
\]
Since $a>1$, this sum extends locally uniformly to
\[
\C\setminus[1,\infty).
\]
Unlike the positive Stieltjes resolvent of Paper III, the signed resolvent is not asserted to be zero-free.

In the smaller domain where the following geometric expansion converges,
\[
\boxed{
\Ccal_{a,\beta,\chi}(z,t)
=
\frac1{1-z}
\sum_{j=0}^{\infty}
\left(-\frac{z}{1-z}\right)^j
\frac1{L(a+\beta j+\ii t,\chi)}.
}
\]

\begin{theorem}[Parameter-shift equation]
The signed resolvent satisfies
\[
\boxed{
(1-z)\Ccal_{a,\beta,\chi}(z,t)
+
z\Ccal_{a+\beta,\beta,\chi}(z,t)
=
\frac1{L(a+\ii t,\chi)}.
}
\]
\end{theorem}

\begin{proof}
Multiply the depth recursion by $z^k$ and sum over $k\ge0$. The left shift contributes
\[
\frac{\Ccal_a-C_a(0,t)}{z}.
\]
Since
\[
C_a(0,t)=L(a+\ii t,\chi)^{-1},
\]
rearrangement gives the result.
\end{proof}

\FigureWide{17_signed_signal}{The signed transfer in logarithmic time.}{The signal is quasiperiodic because every integer packet rotates at its own frequency $\log n$. Increasing the depth changes the active scale but not the frequency law.}

\FigureWide{18_normalized_process}{The unit-average-power process.}{The positive packet energy normalizes the long-time mean power to one. The value at $t=0$ remains a distinguished coherent recurrence.}

\FigureWide{19_time_zero_coherence}{Finite time-zero coherence data.}{The plot is a finite computational record only. It is not used to infer the asymptotic RH bound.}

\FigureWide{20_signed_resolvent}{The signed resolvent.}{The function is analytic on the displayed side of the physical cut. Unlike the positive resolvent, it may have zeros and sign cancellations.}

\FigureWide{21_resolvent_functional_equation}{The transfer-equation residual.}{Three complex evaluation points verify the parameter-shift identity to the truncation accuracy of the number-basis sum.}

\FlowChapter{The Borel--Riesz bridge}

The exponential generating function of the signed depth sequence is
\[
\sum_{k=0}^{\infty}
C_{a,\beta,\chi}(k,t)
\frac{x^k}{k!}
=
\sum_{n\ge1}
\frac{\mu(n)\chi(n)}{n^{a+\ii t}}
\e^{x(1-n^{-\beta})}.
\]
Define
\[
\mathcal R_{a,\beta,\chi}(x,t)
=
x
\sum_{n\ge1}
\frac{\mu(n)\chi(n)}{n^{a+\ii t}}
\e^{-x/n^\beta}.
\]
Then
\[
\boxed{
\sum_{k=0}^{\infty}
C_{a,\beta,\chi}(k,t)
\frac{x^k}{k!}
=
\frac{\e^x}{x}
\mathcal R_{a,\beta,\chi}(x,t).
}
\]

\FigureWide{25_borel_riesz_bridge}{The Borel--Riesz identity.}{The discrete depth sequence and its continuous companion are two exact coordinates of the same M\"obius transfer.}

\FlowChapter{Signed Pad\'e packet models}

Let
\[
c_j=C_{a,\beta,\chi}(j,t)
\]
and define the Hankel determinant
\[
\Delta_m(t)
=
\det[c_{i+j}]_{i,j=0}^{m-1}.
\]

\begin{theorem}[Finite signed Pad\'e reduction]
Assume
\[
\Delta_m(t)\ne0.
\]
Then there is a unique rational function
\[
\boxed{
\Ccal_m(z,t)
=
\frac{P_{m-1}(z,t)}{Q_m(z,t)},
\qquad
Q_m(0,t)=1,
}
\]
such that
\[
\Ccal_{a,\beta,\chi}(z,t)-\Ccal_m(z,t)
=O(z^{2m}).
\]
If the poles of $Q_m^{-1}$ are simple, then
\[
\boxed{
\Ccal_m(z,t)
=
\sum_{r=1}^{m}
\frac{\lambda_r(t)}{1-z\xi_r(t)}
}
\]
and
\[
\boxed{
C_{a,\beta,\chi}(j,t)
=
\sum_{r=1}^{m}
\lambda_r(t)\xi_r(t)^j,
\qquad
0\le j\le2m-1.
}
\]
\end{theorem}

\begin{proof}
Write
\[
Q_m(z)=1+q_1z+\cdots+q_mz^m.
\]
The conditions on coefficients $z^m$ through $z^{2m-1}$ form a linear system whose matrix is a column reversal of the Hankel matrix $[c_{i+j}]$. The determinant assumption gives a unique denominator. The numerator is then the polynomial part of $Q_m\Ccal$ through degree $m-1$. This proves the moment contract. Simple poles give the partial-fraction form.
\end{proof}

\begin{example}[The first signed packet leaves the physical interval]
For the principal case,
\[
\xi_1
=
\frac{c_1}{c_0}
=
-0.5198177546\ldots.
\]
The positive packet spectrum belongs to $[0,1)$. The first signed effective node is already negative. This shows that the signed reduction is an indefinite moment problem rather than a positive one.
\end{example}

\FigureWide{22_pade_nodes}{Positive and signed packet nodes.}{The positive companion reduction stays in the declared spectral interval. The signed M\"obius reduction can leave it because cancellation replaces positivity.}

\FigureWide{23_pade_moment_matching}{Moment-exact finite reduction.}{Four packets preserve the first eight positive or signed moments. The later moments are predictions rather than part of the exact contract.}

\FigureWide{24_hankel_contrast}{Positive and signed Hankel determinants.}{The positive determinants remain positive. The signed determinants change sign and become increasingly ill-conditioned, which exposes the nonbreakdown boundary directly.}

\begin{table}[H]
\centering
\caption{Four-packet Pad\'e nodes and weights for the positive companion and signed M\"obius sequences.}
\small
\begin{tabular}{lrrrr}
\toprule Family & Node real & Node imag. & Weight real & Weight imag. \\ \midrule
positive & 0.000000000 & 0.00e+00 & 1.000000003 & -2.17e-17 \\
positive & 0.750017741 & -6.47e-15 & 0.062527773 & -3.78e-15 \\
positive & 0.890070247 & 9.21e-15 & 0.012629961 & 6.89e-15 \\
positive & 0.973460741 & -1.13e-15 & 0.002770400 & -3.09e-15 \\
signed & -0.000000003 & -0.00e+00 & 0.999999972 & 1.68e-16 \\
signed & 0.750048271 & -1.13e-14 & -0.250305329 & 2.33e-14 \\
signed & 0.890460587 & 1.78e-14 & -0.114859170 & -5.74e-14 \\
signed & 0.965969069 & -2.28e-15 & -0.026908372 & 3.40e-14 \\
\bottomrule
\end{tabular}
\end{table}


\InlinePart{Quantum and recurrence interpretation}

\FlowChapter{The number-basis operator system}

Work in the number basis indexed by $n\ge2$ and define
\[
H\ket n
=
\log n\ket n,
\]
\[
K_\beta\ket n
=
-\log(1-n^{-\beta})\ket n,
\]
and
\[
\Mcal_\chi\ket n
=
\mu(n)\chi(n)\ket n.
\]
These three diagonal operators commute. Let
\[
\ket{\phi_a}
=
\sum_{n\ge2}n^{-a/2}\ket n.
\]
Then the exact operator identity is
\[
\boxed{
C_{a,\beta,\chi}(k,t)
=
\mathbf 1_{\{k=0\}}
+
\bra{\phi_a}
\Mcal_\chi
\e^{-kK_\beta}
\e^{-\ii tH}
\ket{\phi_a}.
}
\]
The isolated term is the $n=1$ packet. It contributes at depth zero and vanishes for every positive depth. Keeping it separate avoids assigning an infinite value to $K_\beta$ at $n=1$. The external quantum construction that motivates this coordinate uses the same logarithmic Hamiltonian and number-basis phase. The additional Euclidean filter and M\"obius sign observable are specific to the present packet system.

\FigureWide{26_quantum_operator}{The operator realization.}{The signed amplitude is a coherent readout. Time averaging produces the positive depth-time kernel and removes only the off-diagonal logarithmic frequencies.}

\FlowChapter{The recurrence process and the RH boundary}

For the principal parameters define
\[
c_k=C_{2,2}(k,0)
\]
and
\[
E_k
=
\Kcal(k,k;0)
=
\sum_{n\ge1}
\frac{\mu(n)^2}{n^4}
\left(1-\frac1{n^2}\right)^{2k}.
\]
The preceding papers prove
\[
E_k
\sim
\frac{3}{2^{5/2}\pi^{3/2}}k^{-3/2}.
\]
Define
\[
\rho_k(t)
=
\frac{C_{2,2}(k,t)}{\sqrt{E_k}}.
\]
The exact time-average identity gives
\[
\boxed{
\lim_{T\to\infty}
\frac1{2T}
\int_{-T}^{T}
\abs{\rho_k(t)}^2\dd t
=1.
}
\]
The time-zero power is
\[
\Gamma_k
=
\abs{\rho_k(0)}^2
=
\frac{\abs{c_k}^2}{E_k}.
\]
The classical sequential criterion and the proved energy asymptotic give
\[
\boxed{
\mathrm{RH}
\iff
\Gamma_k
=O_\varepsilon(k^\varepsilon)
\quad
\text{for every }\varepsilon>0.
}
\]

The positive autocorrelation is therefore fully controlled. The remaining assertion is that the deterministic recurrence time $t=0$ is not polynomially exceptional relative to the unit average power.

\begin{boundary}[The exact remaining obstruction]
Gamma dephasing, scale overlap, character decoherence, positive Gram structure, and finite signed Pad\'e reduction do not prove the principal coherence bound. The missing theorem must use the internal arithmetic signs of $\mu(n)$ inside the principal sector.
\end{boundary}

\FigureWide{27_status_boundary}{The exact boundary.}{The positive depth-time theory is unconditional. Finite signed models are exact under a visible nonbreakdown condition. The principal time-zero coherence remains open.}

\FigureWide{28_function_family_map}{The generated function family.}{The paper adds one positive reproducing kernel, one zero-free Gamma dephasing limit, one signed resolvent, one log-frequency law, one character cross-kernel family, and finite signed Pad\'e packet models.}

\InlinePart{Worked algorithms and verification}

\FlowChapter{Computational route}

\begin{algorithm}[Constructing the depth-time kernel and signed reduction]
\small
\begin{enumerate}[label=\arabic*.,leftmargin=*]
\item Sieve the M\"obius function through the declared cutoff.
\item Form the nodes $y_n=1-n^{-\beta}$.
\item Form the positive weights $\mu(n)^2\abs{\chi(n)}^2n^{-2a}$.
\item Evaluate the depth-time kernel by diagonal contraction.
\item Evaluate the signed transfer with the M\"obius phase retained.
\item Compare the dephased positive kernel with the Gamma-ratio limit.
\item Form signed moments $c_0,\ldots,c_{2m-1}$.
\item Check the Hankel determinant.
\item Solve the Pad\'e denominator system when the determinant is nonzero.
\item Return the $m$ signed packet nodes and weights with a $2m$-moment certificate.
\end{enumerate}
\end{algorithm}

\FlowChapter{Verification ledger}

\[\boxed{18/18\text{ built-in checks passed}}\]


\begin{longtable}{p{0.78\textwidth}p{0.12\textwidth}}
\caption{Built-in computational certificates. Full numerical details are retained in the JSON ledger.}\label{tab:checks}\\
\toprule Check & Result \\ \midrule
\endfirsthead
\toprule Check & Result \\ \midrule
\endhead
depth-time Gram matrix is positive semidefinite & pass \\
finite-difference zeta-ratio coordinates & pass \\
time-average correlation approaches positive kernel & pass \\
principal Gamma modulus closed form & pass \\
log-frequency mean approaches polygamma target & pass \\
log-frequency variance approaches polygamma target & pass \\
signed resolvent parameter-shift equation & pass \\
positive Pade nodes remain in physical interval & pass \\
first signed one-node packet lies outside physical interval & pass \\
positive Pade model matches first eight moments & pass \\
signed Pade model matches first eight moments & pass \\
Borel and Riesz bridge & pass \\
signed finite-difference moments & pass \\
Gamma dephasing asymptotic & pass \\
depth overlap example & pass \\
finite-model long-time mean power & pass \\
character off-diagonal correlation decreases & pass \\
signed resolvent coefficient expansion & pass \\
\bottomrule
\end{longtable}


The independent test suite contains 18 additional tests. It checks the finite-difference formulas, Gram positivity, Gamma limit, hyperbolic-secant overlap, polygamma cumulants, character decoherence, resolvent transfer equation, Pad\'e moment contract, positive versus signed node behavior, Borel--Riesz identity, figure inventory, and result ledger.

\appendix

\chapter{Formula registry}

\begin{longtable}{L{0.30\textwidth}L{0.61\textwidth}}
\toprule
Object & Formula\\
\midrule
\endfirsthead
\toprule
Object & Formula\\
\midrule
\endhead
Signed transfer & $C_{a,\beta,\chi}(k,t)=\sum_n\mu(n)\chi(n)n^{-a-\ii t}(1-n^{-\beta})^k$\\
Depth-time kernel & $\Kcal_{\chi,\psi}(k,\ell;\tau)=\sum_n\mu(n)^2\chi(n)\overline{\psi(n)}n^{-2a-\ii\tau}(1-n^{-\beta})^{k+\ell}$\\
Squarefree ratio & $D_\xi(s)=L(s,\xi)/L(2s,\xi^2)$\\
Gamma shape & $\alpha=(2a-1)/\beta$\\
Gamma dephasing & $\Dcal_{\alpha,\beta}(\tau)=\Gamma(\alpha+\ii\tau/\beta)/\Gamma(\alpha)$\\
Principal modulus & $\abs{\Dcal_{3/2,2}(\tau)}^2=(1+\tau^2)/\cosh(\pi\tau/2)$\\
Depth overlap & $\sech^\alpha(\tfrac12\log(k/\ell))$\\
Log-frequency mean & $-\psi(\alpha)/\beta$\\
Log-frequency variance & $\psi_1(\alpha)/\beta^2$\\
Signed resolvent & $\Ccal(z,t)=\sum_n\mu(n)\chi(n)n^{-a-\ii t}/[1-z(1-n^{-\beta})]$\\
Resolvent transfer & $(1-z)\Ccal_a+z\Ccal_{a+\beta}=L(a+\ii t,\chi)^{-1}$\\
Riesz companion & $\mathcal R(x,t)=x\sum_n\mu(n)\chi(n)n^{-a-\ii t}\e^{-x/n^\beta}$\\
Hankel determinant & $\Delta_m(t)=\det[C(i+j,t)]_{i,j=0}^{m-1}$\\
Pad\'e contract & $C(j,t)=\sum_{r=1}^{m}\lambda_r(t)\xi_r(t)^j$ for $0\le j\le2m-1$\\
Principal coherence & $\Gamma_k=\abs{C_{2,2}(k,0)}^2/\Kcal(k,k;0)$\\
\bottomrule
\end{longtable}

\chapter{Implementation map}

The complete implementation is supplied as
\begin{center}
\texttt{projector\_generated\_special\_functions\_iv.py}.
\end{center}
It contains the M\"obius sieve, character packets, depth-time kernel, finite time averages, exact zeta-ratio checks, Gamma asymptotics, log-frequency laws, signed resolvent, Pad\'e packet reduction, Borel--Riesz transform, figure generation, table generation, and built-in verification ledger.

The central computational path is:
\begin{lstlisting}
params = KernelParameters(a=2.0, beta=2.0, conductor=1)
data = build_number_packet_data(params, nmax=600_000)
kernel = data.positive_kernel(k + ell, tau)
signed = data.signed_transfer(k, t)
numerator, denominator = pade_from_moments(moments, order=m)
nodes, weights = pade_packet_nodes_weights(moments, order=m)
\end{lstlisting}

The positive kernel is evaluated by diagonal contraction. The signed Pad\'e model is built from moments only. Neither path materializes a dense number-basis operator.

\chapter{Complete Python source}

\lstinputlisting{projector_generated_special_functions_iv.py}

\section{Reproduction}

From the source directory, run
\begin{lstlisting}[language=bash,numbers=none]
python projector_generated_special_functions_iv.py
python generate_tables.py
python -m pytest -q
pdflatex -interaction=nonstopmode Projector_Generated_Special_Functions_IV.tex
pdflatex -interaction=nonstopmode Projector_Generated_Special_Functions_IV.tex
pdflatex -interaction=nonstopmode Projector_Generated_Special_Functions_IV.tex
\end{lstlisting}

\section{Source boundary}

The external historical inputs, prior special-function literature, and the 2026 logarithmic-spectrum quantum construction belong in the consolidated end-of-series bibliography. No historical priority claim is made for every displayed specialization. The paper claims the formulas and theorems proved here, the generated function organization, and the computational certificates supplied with the release.

\end{document}
